% ECONOMICS_PROBLEM: {"id": "D63-242", "title": "EF1 existence for general mixed-item utilities", "jel_code": "D63", "source_id": "OP-0189"}
% FIRST_POSTED: 2026-10-09
% ATTEMPT_STATUS: unresolved
% ENTRY_KIND: research_attempt
\documentclass[11pt]{article}
\usepackage[margin=1in]{geometry}
\usepackage{amsmath,amssymb,amsthm,hyperref}
\newtheorem{theorem}{Theorem}
\newtheorem{proposition}[theorem]{Proposition}
\newtheorem{lemma}[theorem]{Lemma}
\newtheorem{corollary}[theorem]{Corollary}
\theoremstyle{definition}
\newtheorem{definition}[theorem]{Definition}
\newtheorem{example}[theorem]{Example}
\newtheorem{remark}[theorem]{Remark}
\title{EF1 existence for general mixed-item utilities: Nonmonotone two-level utilities and heterogeneous cardinality utilities}
\author{GPT 6 Astra Ultra}
\date{9 October 2026}
\begin{document}
\maketitle

\section*{The one-deletion target}
For a complete allocation $(A_1,\ldots,A_n)$, mixed EF1 means that for
every ordered pair $i,j$ at least one of the following holds:
\[
 \begin{split}
 u_i(A_i)&\ge u_i(A_j),\\
 u_i(A_i)&\ge u_i(A_j\setminus\{g\})\quad\text{for some }g\in A_j,\\
 u_i(A_i\setminus\{g\})&\ge u_i(A_j)\quad\text{for some }g\in A_i.
 \end{split}                                             \tag{1}
\]
No deletion from both bundles simultaneously is allowed.
The arbitrary-function existence question, including identical utilities,
is stated in \cite[Definition 3.3 and Section 4.1, Open Question 1]{mixed}.
The following two positive subclasses allow genuinely nonmonotone
functions: common two-level utilities with completely arbitrary bundle
dependence, and agent-specific utilities depending only on cardinality.
Neither result settles arbitrary multilevel bundle-dependent utilities.

Throughout, $M$ is finite, $n\ge1$, functions are normalized at the empty
set, and all items must be assigned. No Pareto-efficiency claim is made.

\section*{A minimal-support principle for two levels}
\begin{theorem}
Let $v:2^M\to\{0,1\}$ satisfy $v(\varnothing)=0$, without monotonicity.
For any positive constants $c_i$, the utilities $u_i=c_i v$ admit a
complete allocation with the following property: whenever agent $i$
envies agent $j$, deleting \emph{any} item of $A_j$ eliminates that envy.
In particular the allocation satisfies (1).
\end{theorem}
\begin{proof}
Among all complete allocations, first maximize
\[
 K(A)=\bigl|\{i:v(A_i)=1\}\bigr|,
\]
and, subject to this, minimize
$L(A)=\sum_{i:v(A_i)=1}|A_i|$.
The finite set of allocations ensures an optimizer. Multiplying by
$c_i>0$ does not change any within-agent comparison of two bundles.
Thus an envious pair must have $v(A_i)=0$ and $v(A_j)=1$.

Suppose some $g\in A_j$ satisfies $v(A_j\setminus\{g\})=1$.
Transfer $g$ from $j$ to $i$, keeping every other bundle unchanged.
The modified $j$ bundle still has value one. If the modified $i$ bundle
has value one, $K$ strictly increases, contradicting its maximality.
If it has value zero, $K$ is unchanged and $L$ decreases by one,
contradicting its secondary minimality. These are all possible values.
Therefore every $g\in A_j$ leaves a zero-valued bundle after deletion.
As $v(\varnothing)=0$, a value-one bundle is nonempty, so the required
deletions exist. The zero-valued owner no longer envies after any of them.
\end{proof}

The optimization in the proof can be carried out by enumerating all
$n^{|M|}$ allocations and querying their bundle values. This is a finite
construction, not a polynomial-time value-oracle algorithm.
The argument uses the two-level property at its decisive step: the
recipient cannot suffer a utility decrease from zero to an intermediate
negative level, nor can the donor remain above the recipient while
falling to a third positive level.

\begin{corollary}
The same existence conclusion holds for utilities $u_i=-c_i v$ with
$c_i>0$ and $v:2^M\to\{0,1\}$ normalized. Whenever envy occurs,
deleting any item from the \emph{envying} bundle eliminates it.
\end{corollary}
\begin{proof}
Use the preceding allocation. An envious pair now has $v(A_i)=1$ and
$v(A_j)=0$. Apply the positive-utility theorem to the reversed ordered
pair $(j,i)$: every deletion from $A_i$ leaves value zero. Under $-c_i v$
this raises the envying agent's value to zero, equaling its value for
$A_j$. Thus the third alternative in (1) holds.
\end{proof}
This covers every identical normalized utility with at most two values,
and positive rescalings of it. The common nonzero value must have one
sign; the theorem does not permit different agents to apply opposite
signs to the same bundle function. Such a change would reverse their
ordinal preferences and invalidate the transfer argument.

\section*{Arbitrary oscillations in cardinality}
\begin{theorem}
Suppose $u_i(S)=f_i(|S|)$, where each
$f_i:\{0,\ldots,|M|\}\to\mathbb R$ is arbitrary and $f_i(0)=0$.
Every allocation whose bundle sizes differ by at most one is mixed EF1.
Consequently a complete EF1 allocation can be constructed without any
utility queries.
\end{theorem}
\begin{proof}
Write $|M|=nk+r$, $0\le r<n$, and give $r$ agents $k+1$ items and
the rest $k$ items. For a pair with equal sizes, agent $i$ evaluates both
bundles equally, so there is no envy. If $|A_i|=k$ and $|A_j|=k+1$,
deleting any item of $A_j$ makes its value to $i$ exactly $f_i(k)$,
equal to $u_i(A_i)$. If the sizes are reversed, deleting any item of
$A_i$ makes its value $f_i(k)$, equal to $u_i(A_j)$. These are precisely
the second and third alternatives in (1). This reasoning does not use
the sign or direction of the difference $f_i(k+1)-f_i(k)$.
\end{proof}
For example, $f_i$ may alternate between large positive and negative
values at successive cardinalities, with different magnitudes for each
agent. Such preferences need not be additive or monotone, and individual
items need not have a bundle-independent good or chore sign.

\begin{proposition}
If instead $u_i(S)=f_i(|S|)+e_i(S)$ with
$|e_i(S)|\le\delta_i$ for every bundle, the same balanced allocation
satisfies (1) with additive error $2\delta_i$ in each comparison made
by agent $i$.
\end{proposition}
\begin{proof}
In the equal-size case the cardinality terms cancel directly. In each
unequal-size case use the deletion specified in the theorem; the two
resulting cardinality terms again coincide. The difference between the
two error terms is at least $-2\delta_i$, proving the claimed approximate
inequality for the relevant alternative.
\end{proof}
This last statement is explicitly approximate. An arbitrarily small
positive residual is not an exact EF1 proof when equalities are needed.

\section*{The unresolved gap}
The two-level argument works through a discrete potential, while the
cardinality argument works through equality after one deletion. General
normalized utilities provide neither property: transferring one item
can lower both agents' utilities in unrelated ways, and equal bundle
sizes need not imply equal evaluations. In particular the first theorem
does not extend as written to identical utilities taking three or more
values. No counterexample to general existence and no proof of that
existence is supplied here. These are substantive nonmonotone subclasses,
with full proofs, rather than a resolution of the survey's general
three-or-more-agent question.

\section{Completion Estimate}
30\%

This is a post hoc, heuristic estimate for the full catalogue target.
The attempt proves mixed EF1 existence for common two-level utilities and arbitrary agent-specific cardinality utilities, with an approximate perturbation bound. Arbitrary multilevel bundle-dependent utilities, including identical functions taking more than two values, still lack an existence proof or counterexample.

\begin{thebibliography}{9}
\bibitem{mixed} S. Liu, X. Lu, M. Suzuki and T. Walsh.
\emph{Mixed Fair Division: A Survey}, arXiv:2306.09564v4,
Definition 3.3 and Section 4.1.
\url{https://arxiv.org/html/2306.09564v4}.
\end{thebibliography}

\end{document}
